\documentclass[12pt,a4paper]{article}
\usepackage[utf8]{inputenc}
\usepackage{xeCJK}
\usepackage{geometry}
\usepackage{enumitem}
\usepackage{titlesec}

% 設定頁面邊距
\geometry{
    top=2.5cm,
    bottom=2.5cm,
    left=2.5cm,
    right=2.5cm
}

% 設定中文字體


\setCJKmainfont{Noto Serif CJK TC}
\setCJKsansfont{Noto Sans CJK TC}
\setCJKmonofont{Noto Sans Mono CJK TC}

% 設定行距
\linespread{1.3}

% 設定標題格式
\titleformat{\section}{\large\bfseries}{\thesection}{1em}{}
\titleformat{\subsection}{\normalsize\bfseries}{\thesubsection}{1em}{}

\begin{document}

\begin{center}
\textbf{\Large 推薦函}
\end{center}

\vspace{0.5cm}

\textbf{致：國立台灣科技大學應用科技所博士班}

\vspace{0.3cm}

\section*{推薦人資訊}
\textbf{姓名：} 謝慕揚醫師\\ \hspace{1cm} \textbf{職稱：} 新竹台大分院教學部副主任\\ \hspace{1cm} \textbf{日期：} 2025年9月26日

\section*{推薦對象}
\textbf{被推薦人：} 黃國良 \hspace{1cm} \textbf{申請項目：} 國立台灣科技大學應用科技所博士班丙組

\section*{推薦內容}

我是新竹台大分院教學部副主任謝慕揚醫師，很榮幸為黃國良同仁撰寫此推薦函。我與黃國良的共事期間始於2024年8月，在這段期間內，我深刻體會到他在專業能力與學術潛力方面的卓越表現。

\subsection*{專業背景與卓越成就}

黃國良於2023年1月16日加入本院教學部，同時開始投入實證醫學(EBM)的推廣工作。在將近兩年的時間內，他在實證醫學與醫學資訊領域展現出非凡的才能，獲得以下重要成就：

\textbf{競賽獲獎表現：}
\begin{itemize}[leftmargin=*]
    \item \textbf{2024年（113年）} 榮獲衛生福利部實證醫學競賽優選
    \item \textbf{2025年（114年）} 榮獲衛生福利部實證醫學競賽第三名
\end{itemize}

\textbf{國際學術影響力：}黃國良的研究成果已獲得國際學術界的認可，充分展現其獨立研究的卓越品質與國際級的學術影響力。2024年，他代表新竹臺大分院參與美國醫學圖書館學會年會（MLA），成功發表關於教育科技的海報：《Rural hospital library navigation system design concept》。此項具前瞻性的研究獲得醫學圖書館專業期刊（Journal of Hospital Librarianship）編輯的高度青睞，並於2025年在該期刊的科技專欄中刊登，有力證明其研究成果能直接影響國際實務領域。

\textbf{廣泛的學術產出：}黃國良是一位高度自發性的研究者，持續在多個醫學資訊相關領域進行深入探索：
\begin{itemize}[leftmargin=*]
    \item 積極參與長庚醫學週學術研討會及實證醫學學會研討會，發表多篇關於實證醫學及教育主題的海報
    \item 2025年於數位學習與教育科技國際研討會發表《地區醫院圖書館電子書資訊服務系統建立與服務創新》海報
    \item 2025年於歐洲醫學教育學年會（AMEE）發表海報：《From preparation to performance: A comparative analysis of team success on evidence-based medicine contest》
\end{itemize}

\textbf{專業認證成就：}憑藉其在2024至2025年間持續發表有關醫學教育、醫學圖書館資訊及實證醫學等議題的學術著作，黃國良已順利取得臺灣醫學圖書館學會認證的高級醫學圖書館管理師資格。這項專業認證不僅是對其專業知識的極高肯定，更突顯了其能將尖端學術研究應用於實務場域的綜合能力。

\subsection*{工作表現與能力展現}

在教學部的工作中，黃國良主要負責實證醫學的推廣任務。他展現出卓越的系統性思維，將複雜的推廣工作有條理地分為三個核心部分：

\begin{enumerate}[leftmargin=*]
    \item \textbf{師資組成} - 建立完整的教學師資體系
    \item \textbf{報名推廣} - 有效推廣並招募學員參與
    \item \textbf{學員組隊訓練} - 系統化的訓練與輔導機制
\end{enumerate}

透過他的努力，本院已成功將實證醫學訓練建立為常態性教學活動，並有效協助醫師、藥師、復健等各科部同仁接受專業訓練，積極參與國內實證醫學競賽。

\subsection*{個人特質與學術潛力}

在與黃國良的共事過程中，我觀察到他具備以下優秀特質：

\textbf{企圖心} - 對於學術研究與專業發展展現強烈的進取精神 \\
\textbf{組織能力} - 能夠系統性地規劃並執行複雜的教學推廣計畫 \\
\textbf{思辨能力} - 在活動籌劃與執行過程中展現優秀的邏輯分析與問題解決能力

這些特質在他所負責的各項活動籌劃與執行過程中都得到充分印證，顯示他具備完成博士學位所需的核心能力。

\subsection*{推薦意見}

基於以上觀察，我認為黃國良同仁的表現充分證明他有能力完成博士班學位。他所展現的專業能力、學術成就，以及在實證醫學推廣工作上的傑出表現，都足以證明他具備在應用科技所博士班深造的學術潛力與研究能力。

我強力推薦黃國良同仁申請貴校應用科技所博士班丙組，相信他必能在貴校的學術環境中發揮所長，為學術研究與專業發展做出重要貢獻。

\vspace{1cm}

\textbf{推薦人簽名：} 謝慕揚醫師 \hspace{1cm} \textbf{職稱：} 新竹台大分院教學部副主任 \hspace{1cm} \textbf{聯絡方式：} G11135@hch.gov.tw

\end{document}